\slajdid{09_1-s028}
\begin{frame}[label=09_1-s028]
\gusttekst\setlength{\abovedisplayskip}{3pt}\setlength{\belowdisplayskip}{3pt}\setlength{\abovedisplayshortskip}{2pt}\setlength{\belowdisplayshortskip}{2pt}
Kod dinamičkog režima sa prelaska $V_{OL}$ na $V_{OH}$ u slučaju da je invertor opterećen kapacitivnošću $C_L$ na izlasku, pošto je u pitanju prosto RC kolo možemo odmah da pišemo
\[\tau=R_LC_L\qquad t_{pLH}=0.69\tau\qquad t_r=2.2\tau\]
Prilikom prelaska $V_{OH}$ na $V_{OL}$ moramo da se poslužimo modelom. Pošto napon na izlazu, zbog kapacitivnosti, ne može trenutno da se promeni, napon $V_{DS}$ će biti visok i tranzistor će raditi u zasićenju. Njegova struja je tada konstantna (zanemarićemo $\lambda$) i jednaka je
\begin{columns}[c,onlytextwidth]\column{.66\textwidth}
\[I_{Dn,D}=\frac{k_n}{2}\frac1{1+\frac{(V_I-V_{Tn})}{L_nE_{Cn}}}(V_I-V_{Tn})^2\]
Pošto posmatramo najgori slučaj, smatraćemo da je $V_I=V_{OHmin}$. Po modelu imamo
\[\tau=C_LR_L\]
\[v(t_0^+)=V_{OH}\]
\[v_0(\infty)=V_{DD}-R_Li_{RC}(\infty)=V_{DD}-R_LI_{Dn}\alert{\ll0}\]
Ostaje pitanje da li je tranzistor bio u zasićenju za vreme ovog procesa.
\column{.32\textwidth}\centering
\begin{tikzpicture}[x=.7cm,y=.53cm,font=\sitneoznake,>=Latex]\draw(0,2.3)--(0,1.85)--(-0.09,1.75)--(0.09,1.65)--(-0.09,1.55)--(0.09,1.45)--(-0.09,1.35)--(0.09,1.25)--(0,1.15)--(0,0.7);\node[right]at(0.13,1.5){$R_L$};\draw(0,2.3)--(0,2.5);\draw(-.2,2.5)--(.2,2.5)node[above left]{$V_{DD}$};\draw(0,.7)--(1.8,.7)node[ocirc]{};\draw(0.9,0.7)--(0.9,-0.015000000000000013);\draw(0.9,-0.135)--(0.9,-0.85);\draw(0.7,-0.015000000000000013)--(1.1,-0.015000000000000013);\draw(0.7,-0.135)--(1.1,-0.135);\node[left]at(0.65,-0.07500000000000001){$C_L$};\draw(.9,-.85)node[ground]{};\draw[->](1.8,-.7)--(1.8,.52);\node[right]at(1.8,.4){$+$};\node[right]at(1.8,-.15){$v_o(t)$};\draw(1.8,-.85)node[ground]{};\draw[->](-.27,2.35)--(-.27,.42)--(.68,.42)--(.68,.02);\node[left]at(-.3,1.6){$i_{RC}(t)$};\node at(2.6,1.5){L u H};\end{tikzpicture}
\par\smallskip\begin{tikzpicture}[x=.7cm,y=.53cm,font=\sitneoznake,>=Latex]\draw(0,2.3)--(0,1.85)--(-0.09,1.75)--(0.09,1.65)--(-0.09,1.55)--(0.09,1.45)--(-0.09,1.35)--(0.09,1.25)--(0,1.15)--(0,0.7);\node[right]at(0.13,1.5){$R_L$};\draw(0,2.3)--(0,2.5);\draw(-.2,2.5)--(.2,2.5)node[above left]{$V_{DD}$};\draw(0,.7)--(1.8,.7)node[ocirc]{};\draw(0.9,0.7)--(0.9,-0.015000000000000013);\draw(0.9,-0.135)--(0.9,-0.85);\draw(0.7,-0.015000000000000013)--(1.1,-0.015000000000000013);\draw(0.7,-0.135)--(1.1,-0.135);\node[left]at(0.65,-0.07500000000000001){$C_L$};\draw(.9,-.85)node[ground]{};\draw[->](1.8,-.7)--(1.8,.52);\node[right]at(1.8,.4){$+$};\node[right]at(1.8,-.15){$v_o(t)$};\draw(1.8,-.85)node[ground]{};\draw(-.7,.7)--(0,.7);\draw(-.7,.7)--(-.7,.2);\draw(-.7,-.08)circle(.28);\draw[->](-.7,.09)--(-.7,-.25);\draw(-.7,-.36)--(-.7,-.65)node[ground]{};\node[left]at(-1.02,-.08){$I_{Dn}$};\draw[->](-.3,1.9)--(-.3,1.05)node[midway,left]{$i_{RC}(t)$};\draw[->](.72,-.02)--(.72,.53)--(.15,.53);\node[left]at(.65,-.7){$i_C(t)$};\node at(2.6,1.5){H u L};\end{tikzpicture}
\end{columns}
\end{frame}
